% TOP_PROBLEM: {"id": 9400068, "title": "The density hypothesis for zeroes of the Riemann zeta function", "category_id": 1, "category": {"id": 1, "name": "number_theory", "display_name": "Number Theory", "description": "Properties of integers, prime numbers, Diophantine equations.", "slug": "number-theory", "order_index": 1, "created_at": "2026-07-31T15:26:25.670Z"}, "status": "partially_solved"}
% FIRST_POSTED: null
% ENTRY_KIND: statement_only
% !TeX program = lualatex
\documentclass[11pt]{article}
\usepackage[margin=25mm]{geometry}
\usepackage{fontspec}
\setmainfont{Latin Modern Roman}
\usepackage{amsmath,amssymb,mathtools}
\usepackage{unicode-math}
\setmathfont{Latin Modern Math}
\usepackage{xurl,hyperref}
\hypersetup{colorlinks=true,urlcolor=blue}
\setcounter{secnumdepth}{0}
\setlength{\parindent}{0pt}
\setlength{\parskip}{5pt}
\setlength{\emergencystretch}{3em}
\newcommand{\R}{\mathbb R}
\newcommand{\N}{\mathbb N}
\newcommand{\Z}{\mathbb Z}
\newcommand{\Q}{\mathbb Q}
\newcommand{\C}{\mathbb C}
\newcommand{\D}{\mathbb D}
\newcommand{\F}{\mathbb F}
\newcommand{\sref}[2]{\hyperlink{source:#1:#2}{[S#2]}}
\newcommand{\cataloguescope}[1]{#1}
\begin{document}
% UnsolvedMath ID 9400068; source catalogue code AMR-093-0068
\setcounter{section}{9400067}
\section{The density hypothesis for zeroes of the Riemann zeta function}\label{9400068}
{\small UnsolvedMath ID 9400068 · Number Theory}
\subsection{Definitions and mathematical statement}

\paragraph{Shared notation.}$\mathbb N$, $\mathbb Z$, $\mathbb Q$, $\mathbb R$, and $\mathbb C$ denote the natural numbers, integers, rationals, reals, and complex numbers, respectively. All additional parameters and hypotheses are specified in the question.


\paragraph{Statement recorded in the supplied source.}
Let $N(\sigma,T)$ count zeros $\rho=\beta+i\gamma$ of the Riemann zeta function with $\beta\geq\sigma$ and $0<\gamma\leq T$, counting multiplicity. For every $1/2\leq\sigma\leq1$ and every $\varepsilon>0$, must
\[N(\sigma,T)\ll_{\sigma,\varepsilon}T^{2(1-\sigma)+\varepsilon}\]
hold as $T\to\infty$?
\par\smallskip\textit{Formulation references:} \sref{9400068}{1}, \sref{9400068}{2}.
\subsection{Short English statement}
Let $N(\sigma,T)$ count zeros $\rho=\beta+i\gamma$ of the Riemann zeta function with $\beta\geq\sigma$ and $0<\gamma\leq T$, counting multiplicity. For every $1/2\leq\sigma\leq1$ and every $\varepsilon>0$, must
\[N(\sigma,T)\ll_{\sigma,\varepsilon}T^{2(1-\sigma)+\varepsilon}\]
hold as $T\to\infty$?
\subsection{Sources}
\begin{itemize}
\item[\hypertarget{source:9400068:1}{S1}] ULAM.AI and the UnsolvedMath Contributors, UnsolvedMath: A Curated Collection of Open Mathematics Problems, record 9400068 (AMR-093-0068). Catalogue attribution under CC BY 4.0.
\url{https://huggingface.co/datasets/ulamai/UnsolvedMath}.
\item[\hypertarget{source:9400068:2}{S2}] Original problem formulation and mathematical context.
\url{https://en.wikipedia.org/w/index.php?title=List_of_unsolved_problems_in_mathematics&oldid=1366636767#Number_theory}.
\end{itemize}
\label{end:9400068}
\end{document}
